\newcommand{\tipoRequerimento}{ausência da parcela DIC/2}
\section{DOS FATOS E DOS DIREITOS}

\begin{enumerate}[resume]

\item O art. 468 da REN 1.000/2021 da Aneel estabelece como deve ser apurado o consumo de energia elétrica destinado à iluminação pública para a situação em que não há medição do ponto de iluminação. 

\Citacao{
    Art. 468. O consumo mensal da energia elétrica destinada à iluminação pública deve ser apurado considerando as seguintes disposições:
    
    I - com medição da distribuidora: nas mesmas condições das demais unidades consumidoras dos Grupos A e B com medição;

    II - com medição amostral da distribuidora: a medição amostral deverá ser extrapolada para os demais pontos de iluminação pública, com o consumo da unidade consumidora que agrega os pontos sendo calculado pelo somatório dos consumos individuais;

    III - com sistema de gestão de iluminação pública do poder público municipal ou distrital: o consumo dos pontos de iluminação abrangidos deve ser apurado a partir das informações do sistema de gestão, observado o art. 26 e demais instruções da ANEEL; e 

    IV - nas demais situações: o consumo mensal por ponto de iluminação deve ser estimado considerando a seguinte expressão:

    \begin{equation}
        Consumo Mensal (KWh) = \frac{Carga x (n \times Tempo - DIC/2)}{1000}
    \end{equation}

    Em que,
    \begin{itemize}
        \item Carga = potência nominal total do ponto de iluminação em Watts, incluídos os equipamentos auxiliares, conforme art. 473, devendo ser proporcionalizada em caso de alteração durante o ciclo.
        \item Tempo = tempo considerado para o faturamento diário da iluminação pública, podendo assumir os seguintes valores:
        \begin{itemize}
            \item 24h - para os logradouros que necessitem de iluminação permanente; ou
            \item Tempo médio anual por município homologado no Anexo I da Resolução Homologatória REH nº 2.590, de 13 de agosto de 2019;
        \end{itemize}
        \item DIC = Duração de Interrupção Individual da unidade consumidora que agrega os pontos de iluminação pública no último mês disponível, conforme cronograma de apuração da distribuidora, em horas, conforme Módulo 8 do PRODIST;
        \item n = número de dias do mês ou o número de dias decorridos desde a instalação ou alteração do ponto de iluminação. (grifou-se)
    \end{itemize}
}
\item Cabe ressaltar que essa metodologia para cálculo do faturamento já estava prevista desde a REN 888/2020 da Aneel.

\item Observe que para o cálculo do consumo por estimativa, deve-se subtrair do tempo total considerado para fins de faturamento, o tempo de Duração de Interrupção Individual da Unidade Consumidora que agrega os pontos de iluminação pública sem medição.

\item Vale ressaltar que a parcela ‘DIC/2’ não tem relação com os descumprimentos de indicadores de continuidade, vale dizer, não deve ser usada apenas quando o DIC registrado for maior do que o limite estabelecido. Ou seja, havendo qualquer interrupção de fornecimento, ainda que não supere o tempo máximo tolerável para efeitos de compensação, ainda assim deve-se haver a parcela do DIC/2 no cálculo do faturamento por estimativa.

\item A Aneel estabeleceu o prazo de 07 de julho de 2021, quase um ano após a REN 888/2020 ter entrado em vigor, para as concessionárias se adequarem aos novos procedimentos de cálculo de faturamento por estimativa:
\Citacao{
    Art. 9º Estabelecer as seguintes datas-limites para as distribuidoras de energia elétrica adequarem os seus procedimentos às alterações promovidas por esta Resolução:

    [...]

    II - 7 de julho de 2021 para adequação ao §4º do art. 21-A, §4º do art. 21-E, §1º do art. 22, inciso IV do art. 24 e §7º do art. 24-A; (grifou-se)
}

\item Ocorre que a Concessionária de Energia ainda não considera a variável DIC/2 no cálculo do faturamento. Basta fazer uma simples verificação no QIP do município para observar que a concessionária não faz qualquer menção à parcela DIC/2.

\item Dessa forma, o faturamento da UC que agrega os pontos de iluminação pública sem medição está sendo a maior desde a data-limite que a concessionária tinha para se adequar aos procedimentos de apuração do consumo.

\end{enumerate}

\section{DOS PEDIDOS}

\begin{enumerate}[resume]
\item Ante o exposto, solicita-se que a Energisa Paraíba: 
    \begin{enumerate}
        \item Devolva os valores cobrados indevidamente, referentes a ausência da parcela ‘DIC/2’ no cálculo do consumo por estimativa.
        \item Efetue a devolução dos valores em dobro acrescidos de juros e correção monetária conforme estabelece o § 2º do art. 323, da Resolução Normativa n° 1.000/2021 da Aneel.
        \item Efetue as devidas correções do memorial de cálculo de faturamento por estimativa, conforme art. 468, da REN 1.000/2021 e Módulo 8 do PRODIST.
        \item Proceda com a devolução do valor integral, por meio de depósito bancário, na conta corrente do Município, conforme estabelece o § 7º, do art. 323 da REN n° 1.000/2021 da Aneel, comunicando por escrito o resultado do deferimento desta reclamação, por meio do endereço eletrônico: \emailEmpresa{}
    \end{enumerate}
\end{enumerate}